\section{Introduction}

A simple trick---randomly shuffling neurons before each training step---outperforms a mathematically guaranteed equivariant architecture when only 100 models are available for training. Yet with 250 models, the structural guarantee wins decisively, and flat-MLP never reaches 90\% of its own peak performance within any of our sampled training sizes up to $N{=}1{,}000$. Why does a symmetry-enforcing architecture fail at the very scale where its inductive bias should matter most?

This question arises in the setting of \emph{weight-space property prediction}: given a collection of trained neural networks (a model zoo), we want to learn an encoder that predicts model properties---such as test accuracy---directly from the network weights. The encoder must handle a fundamental symmetry: permuting the neurons in any hidden layer produces a functionally identical network, yet a different weight vector. Encoders that ignore this symmetry are forced to learn the invariance from data; those that enforce it by construction---\emph{equivariant encoders}---are expected to be more sample-efficient, especially when labeled model zoos are small.

This expectation is intuitive and has motivated a line of powerful equivariant architectures: DWSNets \cite{Navon2023Equivariant}, GNN-NFN \cite{Kofinas2024Graph}, and NFN \cite{Zhou2023Permutation}. These encoders achieve state-of-the-art property prediction on model zoos. However, a critical question has remained unanswered: \emph{How much more sample-efficient are they than plain encoders, and does this advantage hold uniformly across all zoo sizes?}

The difficulty is methodological. Prior equivariant encoder papers use private or custom train/test splits; prior plain encoder papers use different datasets. No controlled, shared-split comparison of equivariant vs.\ plain weight-space encoders at systematically varied training set sizes has been conducted. This means practitioners facing a data constraint---how many models must I train to form a useful zoo?---cannot answer from the existing literature.

A deeper issue lurks beneath the surface. \citet{Dayan2026Expressive} proved that all permutation-equivariant weight-space networks are equivalent in \emph{expressivity} given sufficient data. Their theorem implies that equivariant encoders have no ceiling advantage over plain encoders---the difference must be in \emph{sample complexity}. But the magnitude of this difference, and whether it is uniform across data regimes, is precisely what has not been measured.

We fill this gap with a controlled learning curve study on the ModelZooDataset CIFAR-10 CNN zoo \cite{Schurholt2022Model}, comparing three conditions at training sizes $\{100, 250, 500, 1000, \text{full}\}$: (1) flat-MLP (plain encoder), (2) flat-MLP with permutation augmentation (PermAug, a data-level approximation of equivariance), and (3) GNN-NFN (a structural equivariant encoder). Our key insight: the advantage is real and large, but it is \emph{data-regime dependent}.

Our key findings and contributions are:

\textbf{1. A large sample efficiency advantage.} GNN-NFN reaches 90\% of its peak R$^2$ at $N{\approx}147$ training models; flat-MLP requires the full dataset (${\sim}7{,}000$ models) to reach 90\% of its peak R$^2{=}0.886$, yielding an efficiency ratio of approximately 47$\times$. Even at the largest sampled size ($N{=}1{,}000$), flat-MLP achieves only 84\% of its peak while GNN-NFN exceeds 97\% of its own peak.

\textbf{2. A novel data-regime crossover.} At $N{=}100$, PermAug (R$^2{=}0.138$) outperforms structural equivariance (GNN-NFN R$^2{=}{-}0.016$) (single-seed; multi-seed replication required). At $N{=}250$, the ordering reverses (GNN-NFN 0.767 vs.\ PermAug 0.532 vs.\ flat-MLP 0.449). This crossover---not predicted by prior theory and absent from all prior weight-space learning studies---reveals a minimum-data threshold of approximately 150--250 models for equivariant graph encoders.

\textbf{3. Mechanistic grounding.} We verify permutation equivariance to floating-point precision for GNN-NFN ($\text{max\_diff}{=}1.80{\times}10^{-6}$ across 10,000 checks on trained models) and DWSNets ($\text{max\_diff}{=}7.45{\times}10^{-9}$), confirming the structural basis for the efficiency advantage.

\textbf{4. A reusable evaluation protocol.} Shared-split learning curves at $\{100, 250, 500, 1000, \text{full}\}$ training sizes, 90\%-peak efficiency ratio, bootstrap CI, and equivariance verification form a complete protocol for future encoder comparisons.

The remainder of this paper is organized as follows. Section~\ref{sec:related} situates our work within three streams of prior research. Section~\ref{sec:method} describes our methodology. Section~\ref{sec:experiments} details the experimental setup. Section~\ref{sec:results} presents results. Section~\ref{sec:discussion} discusses implications and limitations. Section~\ref{sec:conclusion} concludes.
